\chapter{Cài Đặt Hệ Thống}

\section{Công nghệ sử dụng}

Ứng dụng được viết bằng Python 3.11. Streamlit cung cấp giao diện và session
state; NetworkX biểu diễn graph; OSMnx đọc/ghi dữ liệu OpenStreetMap; Folium và
streamlit-folium hiển thị bản đồ; Shapely xử lý hình học cạnh; scikit-learn được
OSMnx dùng trong tìm nút gần nhất \cite{networkx,streamlit}. Các phiên bản trực
tiếp được khóa trong \texttt{requirements.txt} để hạn chế sai khác môi trường.

\begin{table}[H]
\centering
\caption{Một số thư viện chính trong phiên bản được kiểm tra}
\begin{tabularx}{\linewidth}{p{4.0cm}p{2.2cm}Y}
\toprule
\textbf{Thư viện} & \textbf{Phiên bản} & \textbf{Vai trò} \\
\midrule
Streamlit & 1.50.0 & Giao diện web cục bộ và quản lý phiên \\
OSMnx & 2.0.6 & Nạp graph OSM, ánh xạ tọa độ, tiện ích địa lý \\
NetworkX & 3.5 & Cấu trúc MultiDiGraph và kiểm chứng oracle trong test \\
Folium / streamlit-folium & 0.20.0 / 0.25.3 & Bản đồ, marker, sự kiện click \\
Shapely & 2.2.0 & Chia hình học cạnh dài \\
Pandas / NumPy & 2.3.2 / 2.2.6 & Bảng dữ liệu và tính toán hỗ trợ \\
\bottomrule
\end{tabularx}
\end{table}

\section{Cấu trúc mã nguồn}

\begin{table}[H]
\centering
\caption{Trách nhiệm của các mô-đun chính}
\begin{tabularx}{\linewidth}{p{5.3cm}Y}
\toprule
\textbf{Mô-đun} & \textbf{Trách nhiệm} \\
\midrule
\texttt{app.py} & Điểm vào duy nhất, chuyển chế độ, widget và trình bày kết quả \\
\texttt{src/shortest\_path.py} & Dijkstra tự cài đặt, nhiều đích, hỗ trợ MultiDiGraph \\
\texttt{src/tsp.py} & Held--Karp cho các chế độ kết thúc \\
\texttt{src/heuristics.py} & Nearest-neighbor và 2-opt \\
\texttt{src/route\_optimizer.py} & Lập ma trận, gọi thuật toán và ghép tuyến \\
\texttt{src/map\_service.py} & Nạp/cache graph, kiểm tra metadata và truy vấn mạng phụ \\
\texttt{src/chase.py} & Luật, trạng thái, Arena, tập niềm tin và chính sách máy \\
\texttt{src/chase\_maps.py} & Sân lưới, sân OSM, chia cạnh và sinh kịch bản bằng seed \\
\texttt{src/chase\_ui.py} & Điều khiển ván, bản đồ, replay và xuất JSON \\
\texttt{tests/} & Unit test, oracle test, AppTest và test dữ liệu \\
\bottomrule
\end{tabularx}
\end{table}

\section{Cài đặt Dijkstra}

Phiên bản trong dự án không gọi trực tiếp hàm shortest-path của NetworkX. Heap
chứa bộ ba khoảng cách, thứ tự chèn và nút để tránh yêu cầu so sánh trực tiếp hai
kiểu định danh nút. Với đa cạnh $(u,v)$, thuật toán chọn trọng số nhỏ nhất hợp lệ
giữa các cạnh song song. Đầu vào bị từ chối nếu trọng số là boolean, không phải
số hữu hạn hoặc âm.

\begin{lstlisting}[float,caption={Giả mã Dijkstra nhiều đích},label={lst:dijkstra}]
distance[source] = 0
push(heap, (0, source))
while heap is not empty and targets remain:
    du, u = pop_min(heap)
    if du is stale: continue
    if u is a target: mark u settled
    for each outgoing neighbor v of u:
        weight = minimum valid parallel-edge weight(u, v)
        if du + weight < distance[v]:
            distance[v] = du + weight
            predecessor[v] = u
            push(heap, (distance[v], v))
return distances and reconstructed paths
\end{lstlisting}

Test dùng NetworkX làm oracle trên nhiều graph sinh ngẫu nhiên. Oracle không làm
cho cài đặt đúng một cách tự động, nhưng giúp phát hiện sai khác về khoảng cách
trên nhiều cấu hình hơn các ví dụ viết tay.

\section{Cài đặt Held--Karp và khôi phục tuyến}

Tập điểm được mã hóa bằng bitmask. Khóa bảng động là $(mask,last)$; giá trị lưu
chi phí và nút trước. Cấu hình điểm cuối được xử lý ở bước chọn trạng thái cuối,
không nhân bản ba thuật toán. Trước khi chạy, ma trận được kiểm tra vuông, hữu
hạn tại các cung cần thiết và nhất quán kích thước.

Sau khi khôi phục thứ tự, bộ điều phối lấy đường đi chi tiết đã lưu từ từng lần
Dijkstra. Khi ghép hai đoạn liên tiếp, nút đầu của đoạn sau bị bỏ để không lặp
nút nối. Tổng quãng đường và thời gian được tính từ cạnh thực của tuyến ghép,
không suy ra từ khoảng cách Euclid giữa marker.

\section{Cài đặt Arena và tiền tính khoảng cách}

\texttt{Arena} sao chép rồi đóng băng graph để tránh thay đổi trong ván. Vì sân
giới hạn tối đa 600 nút, hệ thống chạy Dijkstra từ mọi nguồn và lưu cả khoảng
cách lẫn đường đi. Nếu $N=|V|$, chi phí khởi tạo xấp xỉ
$O(N(|E|+|V|)\log |V|)$; đổi lại, truy vấn khoảng cách trong mỗi lượt là $O(1)$
và việc lấy các nút trong ngân sách chỉ quét danh sách nút.

Đây là lựa chọn phù hợp cho sân nhỏ, nhưng không mở rộng trực tiếp sang graph
350 nghìn nút. Nếu dùng sân lớn, cần chuyển sang Dijkstra theo nhu cầu, landmark,
multi-source search hoặc cache có giới hạn.

\section{Cài đặt một lượt chơi}

Hàm \texttt{step} là điểm chuyển trạng thái duy nhất. Trình tự cài đặt như sau:
\begin{enumerate}
  \item kiểm tra ván còn hoạt động và điểm đến nằm trong ngân sách;
  \item tăng lượt, lan truyền tập niềm tin và áp dụng tiết lộ định kỳ;
  \item xét bắt tại vị trí đội máy, sau đó xét thoát;
  \item loại vị trí đội máy khỏi tập niềm tin khi người trốn còn sống;
  \item gọi chính sách chỉ với \texttt{Observation};
  \item di chuyển đồng thời các đội, xét bắt rồi xét hết lượt;
  \item trả về một \texttt{State} mới thay vì sửa trạng thái cũ.
\end{enumerate}

Tính bất biến giúp replay đơn giản: lịch sử chỉ là danh sách các snapshot. JSON
được xuất với schema version 1, nguồn sân, seed, chiến thuật và toàn bộ trạng
thái cần thiết để kiểm tra diễn biến.

\section{Tích hợp giao diện và session state}

Hai nút \textit{Lộ trình/Trốn tìm} cập nhật \texttt{workspace\_mode} và query
parameter. Nhánh trò chơi trả về trước khi graph định tuyến lớn được nạp, nhờ đó
chế độ này vẫn mở nhanh. Các khóa widget định tuyến được giữ sống trong session
state khi widget tạm ẩn; khi quay lại, tiêu chí, kiểu điểm cuối, lựa chọn điểm và
tùy chọn click bản đồ vẫn còn.

Click bản đồ trò chơi sử dụng đối tượng marker trả về bởi streamlit-folium thay
vì click nền chung. Chỉ các chấm xanh trong phạm vi di chuyển được chấp nhận;
dropdown cung cấp phương án truy cập thay thế và giúp kiểm thử ổn định.

\section{Chuẩn bị và bảo toàn dữ liệu}
\enlargethispage{5\baselineskip}

Graph lớn được quản lý bằng Git LFS. Tệp Pickle giúp nạp nhanh; GraphML là phương
án dự phòng. Metadata chứa bbox, quy mô, nguồn và thời điểm tạo. Script PBF đọc
dữ liệu Việt Nam cục bộ, cắt theo bbox, kiểm tra kết nối các hành lang quan trọng
rồi mới thay gói hiện tại. Thiết kế này tránh phụ thuộc Overpass trong mỗi lần
người dùng tối ưu.

Script chuẩn bị sân trò chơi dùng tệp tạm và chỉ thay tệp đích sau khi graph mới
vượt kiểm tra. Test mô phỏng lỗi tải để bảo đảm tệp sân đang hoạt động không bị
ghi đè bởi dữ liệu dở dang.

\section{Trích đoạn mã nguồn cốt lõi}
\label{sec:trich-doan-ma}

Ba đoạn dưới đây là mã thật trong dự án, không phải giả mã. Chúng được chọn vì mỗi
đoạn tương ứng với một phát biểu lý thuyết ở Chương~2 và là chỗ dễ sai nhất khi
cài đặt.

\subsection{Vòng lặp quy hoạch động Held--Karp}

Đoạn sau là phần truy hồi và ba cách chọn trạng thái cuối trong
\texttt{src/tsp.py}. Điểm đáng chú ý là ba kiểu kết thúc chỉ khác nhau ở bước
chọn trạng thái cuối, còn bảng \texttt{dp} dùng chung; nhờ vậy không có ba bản sao
của cùng một thuật toán.

\begin{lstlisting}[caption={Truy hồi Held--Karp và ba kiểu kết thúc},
                   label={lst:heldkarp}]
for mask in range(1, 1 << n):
    for j in range(n):
        if not mask & (1 << j) or not isfinite(dp[mask][j]):
            continue
        for k in range(n):
            if mask & (1 << k):
                continue
            next_mask = mask | (1 << k)
            candidate = dp[mask][j] + costs[destinations[j]][destinations[k]]
            # Transition: noi them mot diem chua ghe vao cuoi tuyen.
            if candidate < dp[next_mask][k]:
                dp[next_mask][k] = candidate
                parent[next_mask][k] = j

full_mask = (1 << n) - 1
if return_to_start:
    # Closed TSP: chon diem cuoi tam thoi roi cong canh quay ve start.
    last = min(range(n),
               key=lambda j: dp[full_mask][j] + costs[destinations[j]][start_index])
    optimal_cost = dp[full_mask][last] + costs[destinations[last]][start_index]
elif end_index is not None:
    # Fixed end: trang thai cuoi bat buoc nam tai end_index.
    last = destinations.index(end_index)
    optimal_cost = dp[full_mask][last]
else:
    # Free end: lay min tren moi diem cuoi, KHONG cong canh ve start.
    last = min(range(n), key=lambda j: dp[full_mask][j])
    optimal_cost = dp[full_mask][last]
\end{lstlisting}

Việc kiểm tra \texttt{isfinite(dp[mask][j])} trước khi mở rộng là bắt buộc trên
graph có hướng: nếu một trạng thái đã không tới được, mọi trạng thái mở rộng từ nó
cũng vô nghĩa, và phép cộng với $+\infty$ sẽ lan giá trị vô hạn vào bảng. Khi
\texttt{optimal\_cost} vẫn vô hạn ở cuối, hệ thống báo lỗi cho người dùng thay vì
trả về một tuyến không hợp lệ.

\subsection{Lan truyền tập niềm tin trong một lượt}

Đoạn sau trong \texttt{src/chase.py} thực hiện đúng công thức lan truyền ở
Chương~2 và hai lần thu hẹp tập niềm tin.

\begin{lstlisting}[caption={Cập nhật tập niềm tin theo lượt},label={lst:belief}]
turn = obs.turn + 1
belief = arena.expand(obs.belief, obs.rules.runner_budget)
revealed = turn % obs.rules.reveal_every == 0
if revealed:
    belief = frozenset({destination})
obs = replace(obs, turn=turn, belief=belief,
              last_seen=destination if revealed else obs.last_seen,
              last_seen_turn=turn if revealed else obs.last_seen_turn)
...
# Surviving the runner phase also rules out all currently occupied hunter nodes.
obs = replace(obs, belief=obs.belief.difference(obs.hunters))
\end{lstlisting}

Dòng cuối là một suy luận chứ không phải một thủ thuật: nếu người trốn còn sống
sau pha di chuyển của mình, thì vị trí hiện tại của người trốn không thể là một nút
đang có đội truy tìm. Đây là thông tin hợp lệ rút ra từ quan sát công khai, nên
chính sách máy được phép dùng. Ngược lại, biến \texttt{state.runner} không bao giờ
xuất hiện trong \texttt{Observation}.

\subsection{Hàm điểm phối hợp và danh sách rút gọn}

\begin{lstlisting}[caption={Chiến thuật phối hợp chặn lối thoát},label={lst:contain}]
def joint_score(moves):
    belief_score = sum(max(1 / (1 + arena.distance(m, p) / budget) for m in moves)
                       for p in belief) / len(belief)
    exit_score = sum(max(1 / (1 + arena.distance(m, e) / budget) for m in moves)
                     for e in observation.exits) / len(observation.exits)
    return belief_score + 0.35 * exit_score

# Include exit-directed moves so the coverage shortlist cannot hide a blockade.
for index, hunter in enumerate(hunters):
    pools[index] = list(dict.fromkeys([
        *pools[index], *(arena.advance(hunter, e, budget) for e in observation.exits)
    ]))
moves = max(product(*pools), key=joint_score)
\end{lstlisting}

Hai chi tiết triển khai đáng nói. Thứ nhất, \texttt{max} bên trong tổng chứ không
phải tổng của mọi đội: một vị trí nghi ngờ chỉ cần \textit{một} đội ở gần là đã
được bao phủ, nên cộng dồn khoảng cách của cả hai đội sẽ thưởng cho việc hai đội
đứng chồng lên nhau. Thứ hai, các nước đi hướng lối thoát được thêm vào danh sách
rút gọn sau khi đã cắt còn tám ứng viên theo coverage; nếu không, một nước chặn
lối thoát tốt có thể bị loại ngay ở bước cắt và hệ số 0,35 sẽ không bao giờ có cơ
hội phát huy tác dụng.
